Here is the precise specialization of the reduction theorem, including
its orbit and orientation data. In the notation of
\cite[\S\S3.1--3.3]{cornutjetchev}, take
\[
 F_{\rm CJ}=\Q,\quad K_{\rm CJ}=F=\Q(\sqrt{Dh_*}),\quad
 B_{\rm CJ}=M_2(\Q),\quad S=\{\ell_0\}.
\]
The field $F$ is imaginary because $D>0$ and $h_*<0$.
It embeds into $M_2(\Q)$ by its regular representation. The source
algebra has no finite ramification, so H.1 holds and the integer
in H.2 is $1+0+1=2$. Moreover $\ell_0\nmid Dh_*N$ and
\[
  (Dh_*/\ell_0)=(D/\ell_0)(h_*/\ell_0)=1\cdot(-1)=-1
\]
by the present support condition and (O.2). Thus H.3 holds in its
inert, rather than ramified, case. The resulting target algebra is
the definite algebra $B$ ramified precisely at $\ell_0$ and infinity.

Take a source Eichler order $R^{\rm src}=R'\cap R''$ of level $N$.
In \cite[\S3.2]{cornutjetchev} the source subgroup is
$H=\widehat{R^{\rm src}}^{\times}$ and the target subgroup is
$H_S=\widehat R^{\times}$. Indeed, away from $\ell_0$ the local
orders are transported by the fixed algebra isomorphisms. At $\ell_0$
the source subgroup is $\operatorname{GL}_2(\Z_{\ell_0})$,
whose determinant is $\Z_{\ell_0}^{\times}$; its inverse image
under the target reduced norm is exactly the unit group of the
maximal division order. The target level is the prime-to-$S$ part
of $N$, which is $N$ itself. This construction allows arbitrary
prime-power exponents in $N$, including its exponent at two.

Both fine conductors are one. To check this assertion using the
definition in \cite[\S3.3]{cornutjetchev}, an embedding of the maximal
order $\mathcal O_F$ into $R'\cap R''$ embeds it into each maximal
order. The intersection of either maximal order with $F$ is an order
containing $\mathcal O_F$, hence is $\mathcal O_F$. Conversely,
if both intersections are $\mathcal O_F$, their intersection with
$R'\cap R''$ is $\mathcal O_F$. The same argument applies on the
target. Thus $c'=c''=c=c_S=(1)$ in that theorem. Its parity
condition at $\ell_0$ is $0-0\equiv v_{\ell_0}(N)=0\pmod2$.
The local lattice pair there has distance zero. Since $F_{\ell_0}$
is the unramified quadratic extension, its maximal order fixes a
unique vertex of the local tree. Consequently
\[
 N_{\ell_0}(0,0,0)=1,\qquad
 |\operatorname{Gal}(F[c]/F[c_S])|=1,
\]
which also follows directly from
\cite[Lemma~2.1(ii)]{cornutjetchev}. The multiplicity $\kappa$ in
\cite[Theorem~3.1]{cornutjetchev} is therefore exactly one.

It is essential to retain the sign fiber in that theorem. At $\ell_0$
the quotient
\[
 F_{\ell_0}^{\times}\backslash B_{\ell_0}^{\times}/R_{\ell_0}^{\times}
       \simeq\Z/2\Z
\]
is measured by the valuation of the reduced norm modulo two:
norms from $F_{\ell_0}^{\times}$ have even valuation, whereas a
division uniformizer has norm valuation one.
By \cite[\S3.6]{cornutjetchev} this sign is precisely the choice of
residue-field embedding of the unramified quadratic order, with
origin fixed by the chosen reference embedding. Theorem~3.1 gives
a bijection onto one such sign, not onto both. Away from $\ell_0$
its map on local orbit data is a bijection, so we may also restrict
to the chosen orientations at the primes of $N$. The resulting
map is an equivariant bijection of the source class-group torsor
onto the target torsor for that full orientation tuple. Corollary~3.2
gives the same bijection after fixing any supersingular pair in the
target.

Choose the reduction place compatibly with the place used in (O.2).
Sections~4.9--4.10 of \cite{cornutjetchev} identify the projection
of this adelic map with geometric reduction of the corresponding
enhanced elliptic curve on $X_0(N)$. Their local formula is the
identity under the chosen identifications away from $\ell_0$, and
at $\ell_0$ replaces a matrix by a division-uniformizer power
whose exponent is the valuation of its determinant. Thus it is
the same sign and the same reduction map as above. Choose $z_0$
to be the image of the actual reference Heegner point $x_D$ in
(O.1). Writing
\[
 \chi_F([a])=\chi_{h_*}(\operatorname{N}_{F/\Q}a),\qquad
 z_a=[az_0],
\]
equivariance gives the exact finite-sum identity
\[
 \sum_{[a]\in\operatorname{Pic}(\mathcal O_F)}
       \chi_F([a])P_M(z_a)
 =\lambda_M\!\left(
       \overline{L\sum_{[a]\in\operatorname{Pic}(\mathcal O_F)}
                         \chi_F([a])\pi(x_D^{[a]})}\right)
 =\lambda_M(\overline{LP_D})\quad\text{in }\Z/2^M\Z.
\]
Here the supersingular pairs and their parametrization values are
defined over $\mathbb F_{\ell_0^2}$, so every term is in the domain
of the same $\lambda_M$. Additivity, specialization, and multiplication
by the fixed integer $L$ commute with this finite sum. If the Artin
action is written with inverses, reindexing $[a]\mapsto[a]^{-1}$
leaves it unchanged since $\chi_F$ is quadratic. Changing the
reference point within its torsor multiplies the sum only by
$\chi_F$ of that reference class.

